% !TEX program = xelatex
% 2026年版面上项目申请书 — 分文件编辑入口
% 编译命令: xelatex 面上_main.tex && bibtex 面上_main && xelatex 面上_main.tex && xelatex 面上_main.tex
%
% 各部分文件:
%   面上_cover.tex             封面页与摘要（提交时取消注释）
%   面上_part1_立项依据.tex    第一部分：立项依据
%   面上_part2_研究内容.tex    第二部分：研究内容
%   面上_part3_研究基础.tex    第三部分：研究基础
%   面上_part4_其他.tex        第四部分：其他需要说明的情况


\documentclass[12pt]{article}

\usepackage{nsfc}

% 设置楷体字体（nsfc.sty 已加载 ctex，此处勿重复加载）
% AutoFakeBold=true 让 XeLaTeX 合成粗体（STKaiti 无原生粗体）
\setCJKmainfont[AutoFakeBold=true]{STKaiti}
\setmainfont{Times New Roman}

% Citation style: author-year in text, numbered list in bibliography
\usepackage[round,sort]{natbib}

% Additional packages
\usepackage{tcolorbox}
\usepackage{float}
\usepackage{multirow}
\usepackage{array}
\usepackage{longtable}
\usepackage{tabularx}
\usepackage{tikz}

% REF environment (for reference lists within sections)
\setlength{\bibsep}{5pt plus 0.3ex}
\renewcommand\bibsection{\section*{{\fontsize{12pt}{22pt}\selectfont\setlength{\baselineskip}{20pt}}}}
\newenvironment{REF}%
    {\begingroup\fontsize{12pt}{22pt}\selectfont\setlength{\baselineskip}{15pt}}
    {\endgroup}

% Custom theorem-like commands
\newcommand{\conjecture}[2]{\noindent\textbf{\underline{猜想#1}： #2}}
\newcommand{\define}[2]{\noindent\textbf{\underline{定义#1}： #2}}
\newcommand{\lemma}[2]{\vspace{2mm}\noindent\textbf{\underline{引理#1}： #2}}
\newcommand{\property}[2]{\noindent\textbf{\underline{性质#1}: #2}}
\newcommand{\proposition}[2]{\noindent\textbf{\underline{命题#1}： #2}}

% Font size override: 13pt body text with 17pt line spacing
\RequirePackage{anyfontsize}
\renewcommand{\normalsize}{\fontsize{13pt}{17pt}\selectfont}

% Utility commands
\newcommand{\header}{\textbf}
\newcommand{\textzhbf}[1]{\textbf{#1}}
\newcommand{\addImg}[2][1.0]{\includegraphics[width=#1\linewidth]{#2}}
\newcommand{\myEmph}[1]{\textbf{\textcolor[rgb]{0,0,0.25}{#1}}}

\graphicspath{{figures/}}


\begin{document}

% % 封面页与中英文摘要（提交时由 面上_main.tex 通过 \input 引入）

\begin{center}\bf\Large 项目名称（提交时注释掉此封面占位块）\\
Title in English
\end{center}

\textbf{摘要}：
% 字数建议 400 字以内。语气坚定，旗帜鲜明，重点突出：讲明现状、意义、
% 研究目标、内容、思路和预期结果。

PLACEHOLDER

\textbf{Abstract}：

PLACEHOLDER

科学问题属性：

\clearpage
   % 封面页与摘要（提交时取消注释）

\setcounter{page}{1}
%%%%%%%%% TITLE
\title{报告正文（2026版）}

\maketitle
\thispagestyle{empty}

{\large 参照以下提纲撰写，要求内容翔实、清晰，层次分明，标题突出。{\bfseries\color{red}申请书正文原则上不超过30页，鼓励简洁表达。}
\textbf{\nsfClr{请勿删除或改动下述提纲标题及括号中的文字。}}}

%%%%%%%%%%%%%%%%%%%%%%%%%%%%%%%%%%%%%%%%%%%%%%%%%
% 第一部分：立项依据
%%%%%%%%%%%%%%%%%%%%%%%%%%%%%%%%%%%%%%%%%%%%%%%%%
\ContentDes{（一）立项依据：}

\noindent\hspace{\parindent}\nsfClr{(为什么要开展此项研究，研究的科学技术价值如何)}

\subsection{研究背景及科学意义}

PLACEHOLDER

\subsection{国内外相关工作}

PLACEHOLDER

{
\bibliographystyle{gbt7714-author-year}
\bibliography{Cmm}
}


%%%%%%%%%%%%%%%%%%%%%%%%%%%%%%%%%%%%%%%%%%%%%%%%%
% 第二部分：研究内容
%%%%%%%%%%%%%%%%%%%%%%%%%%%%%%%%%%%%%%%%%%%%%%%%%
\ContentDes{（二）研究内容：}

\noindent\hspace{\parindent}\nsfClr{(提纲不做限制，请按照研究工作的自身逻辑撰写。应提炼出特色与创新点、年度研究计划)}

\subsection{研究目标}

PLACEHOLDER

\subsection{研究内容}

PLACEHOLDER

\subsection{拟解决的关键科学问题}

PLACEHOLDER

\subsection{特色与创新之处}

PLACEHOLDER

\subsection{年度研究计划及预期研究结果}

\subsubsection{年度研究计划}

PLACEHOLDER

\subsubsection{预期研究成果}

PLACEHOLDER


\input{面上_part3_研究基础}

\input{面上_part4_其他}

\end{document}
